\section{Đánh giá và Kết luận}
\label{sec:danh-gia-ket-luan}

% ============================================================
% MỤC TIÊU CHƯƠNG:
% Phân tích phản biện, tổng kết, và hướng tương lai.
% Đây là chương để thể hiện hiểu biết sâu và tư duy phản biện.
% Nguồn: Slides "Encode symmetries or not?" (trang 66–68)
% ============================================================

\subsection{Thảo luận: Có nên luôn mã hóa tính đối xứng?}
\label{subsec:thao-luan}
% Đây là phần quan trọng nhất của chương — điểm phản biện khoa học
%
% Slides trình bày 3 góc nhìn (trang 66–68):
%
% Góc nhìn 1 — Hiệu suất tính toán:
%   - Inductive bias giúp với data nhỏ/đắt
%   - Specialized models có thể không cần thiết ở scale lớn
%   - NHƯNG: equivariance vẫn cải thiện scaling laws
%     (Brehmer et al 2025: cùng compute, equivariant model tốt hơn)
%
% Góc nhìn 2 — Loại đối xứng:
%   - Approximate/data-dependent symmetries (eigenvectors, parameter symmetries)
%     vs exact symmetries (rotation, permutation)
%   - Relevance: symmetry có thực sự tồn tại trong task không?
%   - Benchmark issue (Lawrence et al 2025):
%     MỘT SỐ BENCHMARK ĐÃ ĐƯỢC CANONICALIZED SẴN
%     → Thêm symmetry vào có thể GIẢM performance!
%     → Đây là điểm phản biện quan trọng cho nhóm trình bày
%
% Góc nhìn 3 — Flexibility vs Specialization:
%   - Symmetry có thể restrictive cho generative models
%     → Cần symmetry breaking
%   - Foundation models cần flexible/adaptive symmetries
%   - Specialized models vẫn hữu ích như "tools"

\subsection{Ưu điểm và Hạn chế của Geometric Machine Learning}
\label{subsec:uu-nhuoc-diem}
% Ưu điểm:
%   - Data efficiency (chứng minh lý thuyết ở Ch4, thực nghiệm ở Ch5)
%   - OOD generalization và robustness
%   - Giảm số lượng tham số cần học
%   - Universal approximation được bảo toàn (group averaging)
%
% Hạn chế:
%   - Toán học phức tạp, khó thiết kế cho nhóm đối xứng mới
%   - Chi phí tính toán (full group averaging cho large groups)
%   - Symmetry breaking: không tạo được output ít symmetric hơn input
%   - Benchmark có thể đã canonicalized → kết quả misleading

\subsection{Các hướng nghiên cứu tương lai}
\label{subsec:huong-tuong-lai}
% Tổng hợp Open Directions từ slides:
%
%   1. Flexible & adaptive symmetries (không hardcode)
%   2. Any-dimensional models cho PDEs, biology, language
%   3. Topological deep learning (MCN, SMCN)
%   4. Optimal coverings cho groups (symmetry testing)
%   5. Weight space learning ở scale (>2M models)
%   6. Integrate symmetry testing vào architecture design
%   7. Geometry beyond symmetry (Ricci curvature, topo features)

\subsection{Kết luận}
\label{subsec:ket-luan}
% Tóm tắt 3–4 câu:
% - Geometric ML cung cấp framework thống nhất dựa trên symmetry
% - Hai hướng (off-the-shelf và constructed) với trade-offs rõ ràng
% - Lý thuyết đảm bảo lợi ích sample complexity, Expanders giải quyết chi phí
% - Hướng mới: parameter symmetries, weight space learning, flexible symmetries

\subsection{Phân công nhiệm vụ và Tự đánh giá}
\label{subsec:phan-cong}
% Bảng LaTeX 6 cột:
% Họ tên | MSSV | Nhiệm vụ chính | Chương phụ trách | Mức độ hoàn thành (%) | Điểm tự đánh giá
%
% Phân công gợi ý dựa trên cấu trúc mới:
%   - Chương 1 + 2: 1 người (nền tảng lý thuyết)
%   - Chương 3 Hướng 1 + Paper 3 (Phúc): Phúc
%   - Chương 3 Hướng 2 + Paper 2 (Thịnh): Thịnh
%   - Chương 4: 1 người (lý thuyết nâng cao)
%   - Chương 5 (code + thực nghiệm): 1 người
%   - Chương 6 + tổng hợp slide: 1 người

\clearpage

\section*{Phụ lục: Phân công công việc nhóm}
\addcontentsline{toc}{section}{Phụ lục: Phân công công việc nhóm}
% Bảng chi tiết: MSSV | Họ và tên | Phần đóng góp cụ thể

\clearpage
